\begin{table}[htbp]
\centering
\caption{Aggregate annualized decomposition, 2017-Q4--2022-Q4}
\label{tab:latam_regional_decomposition}
\small
\resizebox{\textwidth}{!}{%
\begin{tabular}{lrrrrrr}
\toprule
Aggregation & Income 2017 & Income 2022 & Total growth & Economy composition & Education & Within-group income \\
& \multicolumn{2}{c}{2017 PPP US dollars} & \multicolumn{4}{c}{Percentage points per year} \\
\midrule
Fixed employment weights & 849 & 844 & -0.12 & 0.00 & 1.35 & -1.47 \\
Observed employment weights & 849 & 842 & -0.15 & -0.03 & 1.35 & -1.47 \\
Equal weights by economy & 927 & 872 & -1.20 & 0.00 & 1.00 & -2.21 \\
\bottomrule
\end{tabular}
}
\begin{minipage}{0.98\textwidth}
\footnotesize
\textit{Note:} Fixed employment weights are proportional to each economy's average number of workers across the two endpoints and constitute the baseline. Observed employment weights use each economy's endpoint-specific worker share and therefore include an economy-composition component. Equal weights describe the average economy. Components are common-rescaling allocations and add to total compound annual growth before rounding. Colombia, Ecuador, Paraguay, and Uruguay cross a Tableau series or survey change.
\end{minipage}
\end{table}
